% Macros and theorem environments shared by the web and print versions.
% Theorem-like environments below appear in the dependency graph; numbered within chapters.
\newtheorem{theorem}{Theorem}[chapter]
\newtheorem{proposition}[theorem]{Proposition}
\newtheorem{lemma}[theorem]{Lemma}
\newtheorem{corollary}[theorem]{Corollary}

\theoremstyle{definition}
\newtheorem{definition}[theorem]{Definition}
\newtheorem{example}[theorem]{Example}

\theoremstyle{remark}
\newtheorem{remark}[theorem]{Remark}

% Notation (Keller's conventions: cochain complexes, d of degree +1).
\newcommand{\Ch}{\mathcal{C}}                 % C(B): complexes of B-modules
\newcommand{\Chdg}{\mathcal{C}_{\mathrm{dg}}} % C_dg(B): the dg-category of complexes
\newcommand{\Hcat}{\mathcal{H}}               % H(B): complexes up to homotopy
\newcommand{\Dcat}{\mathcal{D}}               % D(A): derived category
\newcommand{\Mod}{\operatorname{Mod}}
\newcommand{\Hom}{\operatorname{Hom}}
\newcommand{\dgcat}{\mathrm{dgcat}_k}
\newcommand{\Zz}{Z^0}
\newcommand{\Hz}{H^0}
\newcommand{\op}{\mathrm{op}}
\newcommand{\cA}{\mathcal{A}}
\newcommand{\cB}{\mathcal{B}}
\newcommand{\cC}{\mathcal{C}}
\newcommand{\cD}{\mathcal{D}}
\newcommand{\qis}{\mathrm{qis}}
